\documentclass{article}

\usepackage{booktabs}
\usepackage{amsmath,amssymb}
\usepackage{graphicx}
\usepackage{array}
\usepackage[margin=1in]{geometry}

\title{Supplementary Material for Texture Resonance Retrieval for Retrieval-Grounded Editable Audio Effect Preset Selection}
\author{Anonymous Authors}
\date{}

\begin{document}
\maketitle

\renewcommand{\thesection}{S\arabic{section}}
\renewcommand{\thetable}{S\arabic{table}}
\renewcommand{\thefigure}{S\arabic{figure}}
\setcounter{section}{0}
\setcounter{table}{0}
\setcounter{figure}{0}

\section{Canonical Evidence Sources and Metric Notes}

This supplementary package is synchronized to the repository's canonical experimental artifacts and the P0 experiment record. The primary objective evidence consists of the Protocol-A split audit, the earlier resolved-audio-grouped Protocol-A stats, and the P0 E2/E3 summaries that share a 204-query test split. We include the canonical listening-study report as complementary subjective evidence. We do not use the legacy comparison report or the deprecated Protocol-B outputs as formal evidence in this revision.

All objective metrics follow the same operational definitions used in the main paper. \textbf{Param. Dist.} is RMSE over flattened numeric parameter leaves; \textbf{Acc@0.1} measures the fraction of numeric leaves within an absolute tolerance of 0.1; \textbf{Recall} measures retrieval over active numeric parameters; \textbf{Cosine} is cosine similarity in the flattened parameter space; and \textbf{Module} is a schema-level Jaccard index over active DSP modules inferred from keys ending with \texttt{On}. Missing numeric keys are treated as zeros, non-numeric leaves are ignored, and we do not apply per-parameter normalization. Accordingly, absolute magnitudes should be interpreted \emph{within} each stated protocol rather than across protocols.

\section{Protocol-A: Split Audit and Benchmark Definition}

The external benchmark drop-in contains 1267 total records. A post-hoc audit of the original song-name-based held-out list (used in earlier drafts) revealed exact cross-split audio reuse, so the manuscript now anchors its main objective evidence to a stricter resolved-audio-grouped split. Table~\ref{tab:protocolA_split_audit} summarizes both definitions.

\begin{table}[htbp]
	\centering
	\caption{Protocol-A split audit summary. ``Near-duplicate'' counts come from the repository's cheap fingerprint audit and therefore indicate a remaining risk surface rather than a formal perceptual-duplicate guarantee.}
	\label{tab:protocolA_split_audit}
	\small
	\begin{tabular}{lcccc}
		\toprule
		\textbf{Split definition} & \textbf{$N_{\text{test}}$} & \textbf{$N_{\text{kb}}$} & \textbf{Shared audio paths} & \textbf{Near-duplicate pairs} \\
		\midrule
		Legacy song-name split & 211 & 1056 & 16 & 48 \\
		Resolved-audio grouped split & 204 & 1063 & 0 & 200 \\
		\bottomrule
	\end{tabular}
\end{table}

The stricter split groups records by resolved local audio path before sampling the held-out set, thereby eliminating exact cross-split audio reuse under the current audit. At the same time, the benchmark remains a pilot setting: the cheap fingerprint still reports many near-duplicate pairs, and the corpus remains guitar-only. The dataset JSON exposes \texttt{SongName}, \texttt{Style}, \texttt{Feature}, \texttt{Parameters}, and \texttt{Vectors}, but it does not include a separate auditable text-source field; accordingly, \texttt{Style} and \texttt{Feature} are the only explicit text-side provenance fields available in-repo. The stricter split contains 204 paired TRR rows but only 203 unique \texttt{query\_name} labels because \texttt{Hard Rock Crunch} appears twice in the dataset; all pairing therefore uses \texttt{query\_idx}.

\section{P0 Protocol-A: Direct Retrieval Comparison}

P0 Protocol-A compares TRR with Wav2Vec, FeatureNN, CLAP, PaSST, and PANNs on a shared 204-query split and a 1,063-record retrieval database. All methods use top-1 retrieval. This is the main baseline matrix used by the revised manuscript. The earlier resolved-audio-grouped Protocol-A table is retained below as split-audit context, not as the only main result table.

\begin{table}[htbp]
	\centering
	\caption{P0 Protocol-A direct retrieval comparison. Lower is better for L2 and Norm.L2; higher is better for the other metrics.}
	\label{tab:retrieval_methods_direct}
	\small
	\resizebox{\linewidth}{!}{%
	\begin{tabular}{lccccccc}
		\toprule
		\textbf{Method} & \textbf{$n$} & \textbf{L2} $\downarrow$ & \textbf{Norm.L2} $\downarrow$ & \textbf{Acc@0.1} $\uparrow$ & \textbf{Recall} $\uparrow$ & \textbf{Cosine} $\uparrow$ & \textbf{Module} $\uparrow$ \\
		\midrule
		TRR (Ours) & 204 & \textbf{8.7287} & \textbf{0.1881} & \textbf{0.5766} & \textbf{0.5404} & \textbf{0.8415} & \textbf{0.8297} \\
		Wav2Vec & 204 & 22.5174 & 0.2400 & 0.5015 & 0.4648 & 0.5642 & 0.7423 \\
		FeatureNN & 204 & 17.9391 & 0.2236 & 0.5264 & 0.4779 & 0.6030 & 0.7544 \\
		CLAP & 204 & 16.1257 & 0.2298 & 0.5389 & 0.4922 & 0.6531 & 0.7837 \\
		PaSST & 204 & 18.2860 & 0.2132 & 0.5422 & 0.5001 & 0.6192 & 0.7960 \\
		PANNs & 204 & 21.0269 & 0.2260 & 0.5090 & 0.4759 & 0.5826 & 0.7606 \\
		\bottomrule
	\end{tabular}
	}
\end{table}

\section{P0 Protocol-A: Statistical Reporting}

The P0 statistical report tests TRR against each baseline on per-query Norm.L2 using Wilcoxon signed-rank tests with Holm correction. The tests consistently reject the null hypothesis, but the effect sizes are small; the correct interpretation is statistically reliable but moderate TRR advantage.

\begin{table}[htbp]
	\centering
	\caption{P0 Protocol-A: TRR versus baseline statistical tests on per-query Norm.L2.}
	\label{tab:protocolA_ci}
	\small
	\resizebox{\linewidth}{!}{%
	\begin{tabular}{lccccc}
		\toprule
		\textbf{Pair} & \textbf{$n$} & \textbf{$p$-value} & \textbf{Test} & \textbf{Cohen's $d$} & \textbf{Holm reject} \\
		\midrule
		TRR vs Wav2Vec & 204 & $1.564\times10^{-7}$ & Wilcoxon & -0.2101 & Yes \\
		TRR vs FeatureNN & 204 & $9.523\times10^{-5}$ & Wilcoxon & -0.1380 & Yes \\
		TRR vs CLAP & 204 & $4.476\times10^{-4}$ & Wilcoxon & -0.2189 & Yes \\
		TRR vs PaSST & 204 & $1.254\times10^{-3}$ & Wilcoxon & -0.0980 & Yes \\
		TRR vs PANNs & 204 & $2.122\times10^{-5}$ & Wilcoxon & -0.1528 & Yes \\
		\bottomrule
	\end{tabular}
	}
\end{table}

\section{TRR Mechanism Ablations}

The main P0 table includes one mechanism-relevant control: TRR and Wav2Vec use the same Wav2Vec2 acoustic backbone, but differ in the aggregation operator. Wav2Vec uses a first-order mean-pooled embedding, whereas TRR uses projected frame activations followed by Gram-matrix second-order statistics and L2 normalization. Table~\ref{tab:trr_backbone_control} isolates this comparison from the broader external-baseline matrix.

\begin{table}[htbp]
	\centering
	\caption{Backbone-controlled diagnostic from P0 Protocol-A. TRR and Wav2Vec share the Wav2Vec2 backbone but use different aggregation operators.}
	\label{tab:trr_backbone_control}
	\small
	\resizebox{\linewidth}{!}{%
	\begin{tabular}{lcccccc}
		\toprule
		\textbf{Method} & \textbf{Aggregation} & \textbf{Norm.L2} $\downarrow$ & \textbf{Acc@0.1} $\uparrow$ & \textbf{Recall} $\uparrow$ & \textbf{Cosine} $\uparrow$ & \textbf{Module} $\uparrow$ \\
		\midrule
		TRR & Gram second-order statistics & 0.1881 & 0.5766 & 0.5404 & 0.8415 & 0.8297 \\
		Wav2Vec & Mean pooling & 0.2400 & 0.5015 & 0.4648 & 0.5642 & 0.7423 \\
		\bottomrule
	\end{tabular}
	}
\end{table}

To further test the mechanism rather than only the final leaderboard, we recomputed controlled Wav2Vec2 layer features on the same 204-query P0 split with a fixed random seed and compared same-backbone aggregation variants. This ablation is not a replacement for the user-verified P0 E2 leaderboard because it uses the current local execution path and deterministic projection matrices. Its role is narrower: it checks whether the second-order design choices are directionally supported.

\begin{table}[htbp]
	\centering
	\caption{Controlled TRR mechanism ablation on the P0 split using the current local execution path. Lower is better for Norm.L2; higher is better for all other metrics.}
	\label{tab:trr_mechanism_ablation}
	\small
	\resizebox{\linewidth}{!}{%
	\begin{tabular}{lcccccc}
		\toprule
		\textbf{Variant} & \textbf{Layers} & \textbf{Dim} & \textbf{Norm.L2} $\downarrow$ & \textbf{Acc@0.1} $\uparrow$ & \textbf{Recall} $\uparrow$ & \textbf{Cosine} $\uparrow$ \\
		\midrule
		Mean pooling & 5 & 768 & 0.2582 & 0.6089 & 0.4623 & 0.5895 \\
		Gram, no projection & 5 & 589824 & 0.2501 & 0.6103 & 0.4635 & 0.6092 \\
		Gram, projected & 4 & 4096 & 0.2528 & 0.6017 & 0.4585 & 0.6191 \\
		Gram, projected & 5 & 4096 & \textbf{0.2457} & \textbf{0.6104} & \textbf{0.4654} & 0.6224 \\
		Gram, projected & 6 & 4096 & 0.2529 & 0.5985 & 0.4485 & 0.6317 \\
		Gram, projected & 4+5+6, $d=16$ & 256 & 0.2501 & 0.6052 & 0.4449 & 0.6070 \\
		Gram, projected & 4+5+6, $d=32$ & 1024 & 0.2480 & 0.6022 & 0.4640 & \textbf{0.6398} \\
		Gram, projected & 4+5+6, $d=64$ & 4096 & 0.2546 & 0.5963 & 0.4487 & 0.6249 \\
		Gram, projected & 4+5+6, $d=128$ & 16384 & 0.2576 & 0.6054 & 0.4538 & 0.6121 \\
		Gram, no L2 norm & 4+5+6, $d=64$ & 4096 & 0.2849 & 0.5645 & 0.3944 & 0.5857 \\
		\bottomrule
	\end{tabular}
	}
\end{table}

The mechanism result supports three bounded claims. First, Gram aggregation is not merely a Wav2Vec2-backbone effect: the same-layer Gram variants improve normalized distance over same-backbone mean pooling. Second, normalization is important; removing L2 normalization produces the weakest row in the table. Third, architectural choices remain meaningful: layer~5 and $d=32$/$d=64$ projected multi-layer variants trade off different metrics, so the paper should not claim that the reported TRR configuration is globally optimal. The generated artifact also records top-3 and top-5 oracle-best normalized L2 within each retrieval list, which we use only as a top-$K$ sensitivity diagnostic rather than as the deployed top-1 metric. We therefore frame TRR as a useful texture-aware retrieval prior rather than as a fully optimized representation family.

\section{Resolved-Audio-Grouped Protocol-A Audit Context}

Before the P0 baseline matrix was added, the revision used a resolved-audio-grouped Protocol-A table to avoid exact cross-split audio reuse. That table remains useful as audit context because it documents the anti-leakage split behavior and the earlier CLAP-overlap comparison. It should not be mixed numerically with the P0 table because the reported L2 scale and baseline set differ.

\begin{table}[htbp]
	\centering
	\caption{Resolved-audio-grouped Protocol-A audit table retained for split transparency.}
	\label{tab:protocolA_sig}
	\small
	\resizebox{\linewidth}{!}{%
	\begin{tabular}{lcccccc}
		\toprule
		\textbf{Method} & \textbf{$n$} & \textbf{Param. Dist.} & \textbf{Acc@0.1} & \textbf{Recall} & \textbf{Cosine} & \textbf{Module} \\
		\midrule
		TRR (Ours) & 204 & 8.0467 & 0.5169 & 0.4595 & 0.8247 & 0.8051 \\
		Text-RAG & 204 & 24.1705 & 0.4717 & 0.3967 & 0.4943 & 0.7456 \\
		Wav2Vec-RAG & 204 & 23.8197 & 0.4251 & 0.3505 & 0.5292 & 0.7043 \\
		FeatureNN-RAG & 204 & 19.2613 & 0.4198 & 0.3600 & 0.5479 & 0.7134 \\
		CLAP & 201 & 22.3102 & 0.4492 & 0.3971 & 0.5725 & 0.7240 \\
		\bottomrule
	\end{tabular}
	}
\end{table}

\section{Protocol-A: Per-Query Diagnostics}

To make the main benchmark less dependent on mean-only reporting, we derive per-query diagnostics from the canonical Protocol-A CSV for the resolved-audio-grouped split. The source artifact is the derived Protocol-A analysis file pair in Markdown and JSON form.

These diagnostics are paired by \texttt{query\_idx}. The TRR rows contain 204 paired queries but only 203 unique \texttt{query\_name} labels because one label (\texttt{Hard Rock Crunch}) is reused in the CSV; we therefore treat \texttt{query\_idx} as the authoritative pairing key.

\begin{table}[htbp]
	\centering
	\caption{Per-query win/loss/tie summary for TRR under Protocol-A. The interquartile range collapses to zero in many rows because several queries retrieve the same preset under both methods.}
	\label{tab:protocolA_pairwise_diagnostics}
	\small
	\resizebox{\linewidth}{!}{%
	\begin{tabular}{llcccc}
		\toprule
		\textbf{Baseline} & \textbf{Metric} & \textbf{Wins} & \textbf{Losses} & \textbf{Ties} & \textbf{Mean $\Delta$} \\
		\midrule
		CLAP & Param. Dist. & 130 & 54 & 17 & 14.1460 \\
		CLAP & Cosine & 124 & 60 & 17 & 0.2503 \\
		CLAP & Module & 79 & 39 & 83 & 0.0781 \\
		FeatureNN-RAG & Param. Dist. & 130 & 49 & 25 & 11.2146 \\
		FeatureNN-RAG & Cosine & 122 & 57 & 25 & 0.2769 \\
		FeatureNN-RAG & Module & 86 & 40 & 78 & 0.0917 \\
		Text-RAG & Param. Dist. & 117 & 69 & 18 & 16.1238 \\
		Text-RAG & Cosine & 117 & 69 & 18 & 0.3304 \\
		Text-RAG & Module & 73 & 56 & 75 & 0.0595 \\
		Wav2Vec-RAG & Param. Dist. & 125 & 40 & 39 & 15.7730 \\
		Wav2Vec-RAG & Cosine & 121 & 44 & 39 & 0.2956 \\
		Wav2Vec-RAG & Module & 84 & 43 & 77 & 0.1008 \\
		\bottomrule
	\end{tabular}
	}
\end{table}

\begin{table}[htbp]
	\centering
	\caption{Representative Protocol-A L2 cases from the derived per-query analysis. Positive $\Delta$ indicates TRR is better.}
	\label{tab:protocolA_cases}
	\small
	\resizebox{\linewidth}{!}{%
		\begin{tabular}{lll>{\raggedright\arraybackslash}p{3.5cm}>{\raggedright\arraybackslash}p{3.5cm}c}
		\toprule
		\textbf{Type} & \textbf{Baseline} & \textbf{Query} & \textbf{TRR retrieved} & \textbf{Baseline retrieved} & \textbf{$\Delta$L2} \\
		\midrule
		Largest gain & CLAP & Modular Hybrid Drift & Prepared Guitar Static & Ethereal Wave & +71.6821 \\
		Largest gain & FeatureNN-RAG & Sub-Bass Drone & Hauntology Library & Dark Ambient Drone Prism & +72.1369 \\
		Largest gain & Text-RAG & Sub-Bass Drone & Hauntology Library & Studio Clean & +72.2040 \\
		Largest gain & Wav2Vec-RAG & 3-Chord Punk & Distorted Cassette & Britpop Jangle & +76.2671 \\
		Largest failure & CLAP & Talk Box Funk & Schlager Pop & Drum and Bass Synth & -57.3198 \\
		Largest failure & FeatureNN-RAG & Rave Hardcore & Dark Ambient & Melodic Death Metal & -56.3859 \\
		Largest failure & Text-RAG & Talk Box Funk & Schlager Pop & Studio Clean & -57.3623 \\
		Largest failure & Wav2Vec-RAG & Phased Synth & London Jazz & Nightcore Speed & -51.6878 \\
		\bottomrule
	\end{tabular}
	}
\end{table}

\section{P0 Protocol-C: Real-Degradation Robustness}

P0 Protocol-C evaluates realistic audio degradations by corrupting raw audio first and then re-encoding the degraded signal into TRR embeddings. It reuses the P0 Protocol-A 204-query test split, uses \texttt{Style}+\texttt{Feature} as the text branch, and computes ranking relevance from parameter-space oracle neighbors rather than exact song-name matching.

\begin{table}[htbp]
	\centering
	\caption{P0 Protocol-C overall summary on the shared 204-query split. Lower is better for L2; higher is better for all other metrics.}
	\label{tab:protocolC_summary}
	\small
	\resizebox{\linewidth}{!}{%
	\begin{tabular}{lccccc}
		\toprule
		\textbf{Method} & $\boldsymbol{\alpha}$ & \textbf{L2} $\downarrow$ & \textbf{Acc@0.1} $\uparrow$ & \textbf{MRR} $\uparrow$ & \textbf{NDCG@5} $\uparrow$ \\
		\midrule
		fixed fusion 0.30 & 0.3000 & 29.2322 & 0.4846 & 0.4095 & 0.3723 \\
		fixed fusion 0.50 & 0.5000 & 20.5551 & 0.6085 & 0.5607 & 0.5528 \\
		fixed fusion 0.70 & 0.7000 & 20.1868 & 0.6104 & 0.5652 & 0.5551 \\
		adaptive fusion & 0.5980 & \textbf{20.1491} & \textbf{0.6121} & \textbf{0.5654} & \textbf{0.5566} \\
		\bottomrule
	\end{tabular}
	}
\end{table}

\begin{table}[htbp]
	\centering
	\caption{P0 Protocol-C per-condition winners across the main multimodal comparison set.}
	\label{tab:protocolC_l2tests}
	\small
	\resizebox{\linewidth}{!}{%
	\begin{tabular}{lcccc}
		\toprule
		\textbf{Condition} & \textbf{Best L2} & \textbf{Best Acc@0.1} & \textbf{Best MRR} & \textbf{Best NDCG@5} \\
		\midrule
		AWGN 20 dB & fixed 0.70 & adaptive & adaptive & adaptive \\
		AWGN 10 dB & fixed 0.70 & adaptive & adaptive / fixed 0.50 & fixed 0.50 \\
		AWGN 5 dB & fixed 0.70 & fixed 0.70 & fixed 0.70 & fixed 0.70 \\
		MP3 128k & fixed 0.70 & fixed 0.50 & fixed 0.50 & fixed 0.50 \\
		MP3 64k & fixed 0.70 & fixed 0.50 & fixed 0.50 & fixed 0.50 \\
		MP3 32k & adaptive & fixed 0.50 & adaptive / fixed 0.50 & fixed 0.50 \\
		Reverb 0.6 & adaptive & fixed 0.50 & adaptive / fixed 0.50 & adaptive \\
		Reverb 1.0 & adaptive & fixed 0.50 & adaptive / fixed 0.50 & fixed 0.50 \\
		Truncate 50\% & adaptive & adaptive & adaptive / fixed 0.70 & fixed 0.70 \\
		Truncate 30\% & adaptive / fixed 0.70 & adaptive & adaptive / fixed 0.70 & adaptive / fixed 0.70 \\
		\bottomrule
	\end{tabular}
	}
\end{table}

The per-condition view shows that adaptive fusion is not a trivial uniform winner; fixed 0.50 and fixed 0.70 remain competitive on several conditions. The stronger conclusion is that audio-heavy fusion is fragile, whereas adaptive fusion consistently moves the system toward the stable mid-to-high text-weight region while preserving multimodal behavior.

\paragraph{Local boundary audit.}
We additionally ran a Protocol-C boundary script with explicit audio-only ($\alpha=0.0$), pure-text ($\alpha=1.0$), and parameter-oracle upper-bound rows. Because this run follows a separate local execution path from the P0 summary artifact, we report it as an audit boundary rather than as a replacement for Table~\ref{tab:protocolC_summary}. The result in Table~\ref{tab:protocolC_boundary_audit} confirms the main qualitative interpretation: audio-only retrieval is weak under realistic degradation, mid/high text weights are stable, and the pure-text boundary remains a strong practical boundary. The oracle row is not a method; it selects the nearest KB item using ground-truth parameters and only quantifies remaining headroom. Thus, adaptive fusion should be read as a multimodal operating strategy near the stable region, not as proof of superiority over text-only retrieval.

\begin{table}[htbp]
	\centering
	\caption{Current-code Protocol-C boundary audit with audio-only and pure-text rows. This table is an execution-path audit artifact, not a replacement for the canonical P0 server summary.}
	\label{tab:protocolC_boundary_audit}
	\small
	\resizebox{\linewidth}{!}{%
	\begin{tabular}{lccccc}
		\toprule
		\textbf{Method} & $\boldsymbol{\alpha}$ & \textbf{L2} $\downarrow$ & \textbf{Acc@0.1} $\uparrow$ & \textbf{MRR} $\uparrow$ & \textbf{NDCG@5} $\uparrow$ \\
		\midrule
		audio-only & 0.0000 & 102.5577 & 0.4962 & 0.1184 & 0.0731 \\
		fixed fusion 0.30 & 0.3000 & 81.0390 & 0.5980 & 0.3916 & 0.3588 \\
		fixed fusion 0.50 & 0.5000 & 66.9431 & 0.6791 & 0.5633 & 0.5463 \\
		fixed fusion 0.70 & 0.7000 & 63.9030 & 0.6832 & 0.5730 & 0.5528 \\
		adaptive fusion & 0.5980 & 65.4614 & 0.6837 & 0.5719 & 0.5529 \\
		pure text & 1.0000 & \textbf{61.5079} & \textbf{0.6872} & \textbf{0.5847} & \textbf{0.5586} \\
		parameter oracle upper bound & oracle & 1.3284 & 0.9548 & 1.0000 & 1.0000 \\
		\bottomrule
	\end{tabular}
	}
\end{table}

\section{Listening Study Statistics}

The listening study uses a multiple-stimulus protocol with a hidden reference and no explicit low-quality anchor. The raw logs label our system as \texttt{HCAP}; throughout the manuscript and this supplement we refer to the same condition as the \emph{TRR-based system}. Across all trials, the study contains 26 unique participants and 910 ratings.

The current logs preserve repeated-measures structure and participant IDs, but they do not preserve device metadata, participant expertise labels, or the provenance and time budget of the manual baseline. We therefore interpret Trials~2--5 as perceptual context rather than as a fairness-critical human-vs-system optimization benchmark.

\begin{table}[htbp]
	\centering
	\caption{Study-level listening-test summary.}
	\label{tab:listening_overview}
	\small
	\begin{tabular}{lccccc}
		\toprule
		\textbf{Trial block} & \textbf{Rows} & \textbf{Participants} & \textbf{Stimuli} & \textbf{Raw mean} & \textbf{Raw median / std} \\
		\midrule
		Trial 1 & 208 & 26 & 8 & 72.92 & 77.50 / 22.72 \\
		Trials 2--5 & 312 & 26 & 3 & 69.53 & 73.50 / 23.53 \\
		Trials 6--10 & 390 & 26 & 3 & 57.54 & 59.00 / 32.24 \\
		\bottomrule
	\end{tabular}
\end{table}

\begin{table}[htbp]
	\centering
	\caption{Repeated-measures nonparametric tests for the listening study.}
	\label{tab:listening_tests}
	\small
	\resizebox{\linewidth}{!}{%
	\begin{tabular}{llccccc}
		\toprule
		\textbf{Block} & \textbf{Comparison} & \textbf{$n$} & \textbf{Mean(A-B)} & \textbf{Median(A-B)} & \textbf{rbc} & \textbf{$p$ (Holm)} \\
		\midrule
		Trials 2--5 & Friedman omnibus ($Q=41.6154$, $p_{\mathrm{perm}}=5\times10^{-5}$) & N/A & N/A & N/A & N/A & N/A \\
		Trials 2--5 & reference vs TRR-based system & 26 & 13.78 & 8.38 & 0.860 & 0.00015 \\
		Trials 2--5 & reference vs manual & 26 & 33.61 & 26.00 & 0.983 & 0.00015 \\
		Trials 2--5 & TRR-based system vs manual & 26 & 19.83 & 17.25 & 1.000 & 0.00015 \\
		Trials 6--10 & Friedman omnibus ($Q=36.2308$, $p_{\mathrm{perm}}=5\times10^{-5}$) & N/A & N/A & N/A & N/A & N/A \\
		Trials 6--10 & reference vs MusicGen & 26 & 47.72 & 50.90 & 1.000 & 0.00015 \\
		Trials 6--10 & reference vs TRR-based system & 26 & 46.71 & 51.30 & 0.994 & 0.00015 \\
		Trials 6--10 & MusicGen vs TRR-based system & 24 & -1.10 & 0.80 & 0.067 & 0.7797 \\
		\bottomrule
	\end{tabular}
	}
\end{table}

\begin{table}[htbp]
	\centering
	\caption{Trial 1 style-matching stimulus labels used in the listening study.}
	\label{tab:listening_style_labels}
	\small
	\begin{tabular}{ll}
		\toprule
		\textbf{Label} & \textbf{Interpretation} \\
		\midrule
		Jazz Clean & clean jazz tone target \\
		Blues Solo & blues lead tone target \\
		Phase & phaser/phase-effect target \\
		Ambient Guitar & ambient guitar texture target \\
		Flanger & flanger-effect target \\
		Modern Metal & modern high-gain tone target \\
		Chorus & chorus-effect target \\
		\bottomrule
	\end{tabular}
\end{table}

\section{Near-Duplicate Sensitivity Analysis}

To assess the impact of near-duplicate presets on retrieval metrics, we progressively remove near-duplicates from the knowledge base using normalized parameter-space RMSE thresholds and re-evaluate locally available stored-vector retrieval methods under the P0 deterministic split. This is a sensitivity analysis rather than the main P0 leaderboard, because the local JSON used for this run contains CLAP vectors for 166 of the 204 test queries and does not contain local PaSST/PANNs vectors. Table~\ref{tab:near_dup_sensitivity} reports the resulting robustness pattern.

\begin{table}[htbp]
	\centering
	\caption{Near-duplicate sensitivity analysis under normalized parameter-space RMSE deduplication. KB size decreases as the threshold $\tau$ increases. Param.\,Dist.\,$\downarrow$; all other metrics $\uparrow$.}
	\label{tab:near_dup_sensitivity}
	\small
	\resizebox{\linewidth}{!}{%
	\begin{tabular}{lccccccc}
		\toprule
		\textbf{Method} & $\tau$ & \textbf{KB size} & \textbf{Param.\,Dist.} $\downarrow$ & \textbf{Acc@0.1} $\uparrow$ & \textbf{Recall} $\uparrow$ & \textbf{Cosine} $\uparrow$ & \textbf{Module} $\uparrow$ \\
		\midrule
		\multicolumn{8}{l}{\textit{$\tau = 0$ (full KB, baseline)}} \\
		TRR & 0 & 1063 & 37.9848 & 0.6605 & 0.5311 & 0.8149 & 0.8212 \\
		Wav2Vec & 0 & 1063 & 74.2198 & 0.6012 & 0.4570 & 0.5895 & 0.7401 \\
		FeatureNN & 0 & 1063 & 60.3766 & 0.6320 & 0.4861 & 0.6688 & 0.7582 \\
		CLAP & 0 & 1063 & 69.1914 & 0.5654 & 0.3941 & 0.6276 & 0.7271 \\
		\midrule
		\multicolumn{8}{l}{\textit{$\tau = 0.005$}} \\
		TRR & 0.005 & 469 & 36.9790 & 0.6461 & 0.5103 & 0.8169 & 0.7866 \\
		Wav2Vec & 0.005 & 469 & 98.7481 & 0.5616 & 0.4254 & 0.5060 & 0.7125 \\
		FeatureNN & 0.005 & 469 & 86.6753 & 0.5838 & 0.4458 & 0.5414 & 0.7127 \\
		CLAP & 0.005 & 469 & 109.5363 & 0.5036 & 0.3384 & 0.4435 & 0.6619 \\
		\midrule
		\multicolumn{8}{l}{\textit{$\tau = 0.01$}} \\
		TRR & 0.01 & 409 & 35.4043 & 0.6415 & 0.5142 & 0.8442 & 0.8040 \\
		Wav2Vec & 0.01 & 409 & 103.6916 & 0.5562 & 0.4330 & 0.5113 & 0.7263 \\
		FeatureNN & 0.01 & 409 & 96.9060 & 0.5698 & 0.4421 & 0.5180 & 0.7064 \\
		CLAP & 0.01 & 409 & 123.1859 & 0.4854 & 0.3385 & 0.4094 & 0.6633 \\
		\midrule
		\multicolumn{8}{l}{\textit{$\tau = 0.02$}} \\
		TRR & 0.02 & 160 & 47.7463 & 0.6203 & 0.4753 & 0.7758 & 0.7612 \\
		Wav2Vec & 0.02 & 160 & 100.7927 & 0.5352 & 0.3916 & 0.5277 & 0.6986 \\
		FeatureNN & 0.02 & 160 & 102.8913 & 0.5261 & 0.3586 & 0.5082 & 0.6713 \\
		CLAP & 0.02 & 160 & 125.3818 & 0.4887 & 0.3142 & 0.4414 & 0.6547 \\
		\midrule
		\multicolumn{8}{l}{\textit{$\tau = 0.05$}} \\
		TRR & 0.05 & 83 & 51.4764 & 0.6023 & 0.4405 & 0.7024 & 0.7136 \\
		Wav2Vec & 0.05 & 83 & 71.8281 & 0.5425 & 0.3663 & 0.5711 & 0.6719 \\
		FeatureNN & 0.05 & 83 & 60.3433 & 0.5470 & 0.3494 & 0.5722 & 0.6472 \\
		CLAP & 0.05 & 83 & 70.0935 & 0.5346 & 0.3431 & 0.6039 & 0.6840 \\
		\bottomrule
	\end{tabular}
	}
\end{table}

\section{Parameter-Cluster Hard Split}

The main text reports a compact hard-split table. This supplementary section records the exact local audit boundary. Starting from the 204-name resolved-audio Protocol-A query list, we cluster records whose normalized parameter RMSE is at most $\tau=0.02$ and remove every knowledge-base candidate in a selected query cluster. The resulting split keeps 46 hard queries, removes 158 requested query rows as same-cluster conflicts, and leaves 142 candidate presets. This is intentionally severe: it tests duplicate-family sensitivity under a small candidate pool and should not replace the main P0 leaderboard.

\begin{table}[htbp]
	\centering
	\caption{Parameter-cluster hard split at $\tau=0.02$. Rows with zero coverage indicate that the reviewed local dataset does not contain usable cached vectors for that method.}
	\label{tab:supp_hard_split_results}
	\small
	\resizebox{\linewidth}{!}{%
	\begin{tabular}{lcccccccc}
		\toprule
		\textbf{Method} & \textbf{$n$} & \textbf{Coverage} & \textbf{KB indexed} & \textbf{Norm.L2} $\downarrow$ & \textbf{L2} $\downarrow$ & \textbf{Acc@0.1} $\uparrow$ & \textbf{Recall} $\uparrow$ & \textbf{Cosine} $\uparrow$ \\
		\midrule
		TRR & 46 & 1.000 & 142 & 0.2424 & 29.8574 & 0.6988 & 0.5843 & 0.7739 \\
		Wav2Vec & 46 & 1.000 & 137 & 0.2595 & 30.5715 & 0.6840 & 0.5540 & 0.7403 \\
		FeatureNN & 46 & 1.000 & 137 & 0.2634 & 34.4328 & 0.6706 & 0.5659 & 0.7487 \\
		CLAP & 35 & 0.761 & 41 & 0.2065 & 31.8661 & 0.7162 & 0.6233 & 0.8402 \\
		PaSST & 0 & 0.000 & 0 & -- & -- & -- & -- & -- \\
		PANNs & 0 & 0.000 & 0 & -- & -- & -- & -- & -- \\
		\bottomrule
	\end{tabular}
	}
\end{table}

The hard split changes the interpretation of the retrieval results. Among full-coverage stored-vector rows, TRR remains the best normalized-distance method. CLAP is stronger on the covered subset, but it covers only 35 of the 46 hard queries and indexes only 41 of 142 candidate presets in this local artifact. Thus the CLAP row is useful as a warning against overclaiming TRR dominance, not as a fully matched replacement for the P0 CLAP comparison.

\section{Uncached Latency Audit}

The cached latency table in the main text measures retrieval scoring after embeddings are already available. A separate audit of the reviewed artifact set found no verified live path that jointly performs audio preprocessing, Wav2Vec2 forward inference, TRR re-encoding, search, and validation for newly uploaded reference audio. Therefore, the paper uses the cached latency table only as a retrieval-scoring diagnostic and does not claim end-to-end live reference-audio latency.

\begin{table}[htbp]
	\centering
	\caption{Uncached latency evidence status for newly supplied reference audio.}
	\label{tab:uncached_latency_status}
	\small
	\begin{tabular}{ll}
		\toprule
		\textbf{Required component} & \textbf{Evidence status} \\
		\midrule
		Audio preprocessing & Not verified in artifact \\
		Wav2Vec2 forward pass & Not verified in artifact \\
		TRR re-encoding & Not verified in artifact \\
		Nearest-neighbor search & Cached-path timing available \\
		Parameter validation & Not verified in uncached path \\
		\bottomrule
	\end{tabular}
\end{table}

\section{Additional Baseline Status}

P0 closes the PaSST and PANNs retrieval-baseline gap for the shared Protocol-A benchmark. We also add a diagnostic direct-regression baseline using mean-pooled Wav2Vec2 features and a two-layer MLP trained only on the P0 knowledge base. This regressor is useful as a boundary condition, but it is not an executable-preset retriever: it predicts dense continuous parameter leaves and still requires constraint projection before deployment. Learned reranking remains future work.

\begin{table}[htbp]
	\centering
	\caption{P0 Protocol-A additional baseline status.}
	\label{tab:additional_baselines_ci}
	\small
	\resizebox{\linewidth}{!}{%
	\begin{tabular}{lcccc}
		\toprule
		\textbf{Method} & \textbf{$n$} & \textbf{Evidence status} & \textbf{Use in paper} & \textbf{Boundary} \\
		\midrule
		PaSST & 204 & P0 complete & Main baseline matrix & Small TRR effect size \\
		PANNs & 204 & P0 complete & Main baseline matrix & Small TRR effect size \\
		MLP-Regressor & 204 & Diagnostic complete & Boundary baseline & Dense non-executable output \\
		\bottomrule
	\end{tabular}
	}
\end{table}

The diagnostic MLP regressor obtains raw L2 62.0612, normalized L2 0.1607, Acc@0.1 0.7884, Recall 0.3209, and Cosine 0.4964 on 204 test queries before range projection. This mixed profile is informative: the low recall and cosine score indicate that direct regression does not recover the same active parameter structure as retrieval. Under the current E8 diagnostic, the range-projected MLP is weaker than TRR top-1 and EPR-K5 on Norm.L2; its role is therefore boundary evidence for non-executable regression, not a stronger numeric method.

\section{E8 Executable-Neighborhood Diagnostics}

E8 evaluates whether retrieval produces a parameter-transferable executable neighborhood rather than only a single nearest preset. \textbf{PNR@K} is parameter-neighborhood recall at threshold 0.05 over the retrieved set. EPR denotes exemplar-preserving projection: a softmax-weighted blend over the retrieved executable presets followed by range projection, with provenance reported as effective exemplars and maximum exemplar weight. Module-aware reranking is included only as a diagnostic negative result because it worsens Norm.L2 and degrades recall, cosine, and module alignment in the current run.

\begin{table}[htbp]
	\centering
	\caption{E8 executable-neighborhood diagnostics on the Protocol-A query list under the current E8 execution path. These rows are diagnostic and not a replacement for the P0 Protocol-A direct-retrieval table. For pure TRR rows, Norm.L2 and Acc@0.1 are top-1 retrieval metrics; TRR top-K reports only neighborhood recall. Provenance is reported as effective exemplars / max weight.}
	\label{tab:executable_neighborhood_supp}
	\small
	\resizebox{\linewidth}{!}{%
	\begin{tabular}{lcccccc}
		\toprule
		\textbf{Method} & \textbf{K} & \textbf{PNR@K} & \textbf{Norm.L2} & \textbf{Acc@0.1} & \textbf{EditCost} & \textbf{Provenance} \\
		\midrule
		TRR top-1 & 1 & 0.4608 & 0.1454 & 0.7863 & -- & -- \\
		TRR top-K & 5 & 0.7500 & -- & -- & -- & -- \\
		EPR-K3 & 3 & 0.6520 & 0.1356 & 0.7872 & 4.1972 & 2.6857 / 0.4664 \\
		EPR-K5 & 5 & 0.7500 & 0.1334 & 0.7908 & 4.1359 & 4.2129 / 0.3438 \\
		MLP + RangeProjection & -- & -- & 0.1603 & 0.7894 & -- & none \\
		Module-aware rerank & 5 & -- & 0.1479 & 0.7870 & -- & diagnostic negative \\
		\bottomrule
	\end{tabular}
	}
\end{table}

The primary hybrid row is EPR-K5, not the range-projected MLP. EPR-K5 gives the best normalized distance in this diagnostic table while retaining exemplar provenance. The range-projected MLP remains useful as a direct-regression boundary row, but it has no retrieved-exemplar provenance.

\section{Trials 2--5: Detailed Per-Trial Results}

The main text reports only the aggregate Trials~2--5 results. Table~\ref{tab:trials2to5_detailed} provides per-trial stimulus labels and per-condition ratings for completeness.

\begin{table}[htbp]
	\centering
	\caption{Trials 2--5 per-trial breakdown. Ratings are raw means $\pm$ standard deviation across 26 participants.}
	\label{tab:trials2to5_detailed}
	\small
	\begin{tabular}{lcccc}
		\toprule
		\textbf{Trial} & \textbf{Stimulus} & \textbf{Reference} & \textbf{TRR-based} & \textbf{Manual} \\
		\midrule
		Trial 2 & Jazz Clean & $91.19 \pm 10.52$ & $75.12 \pm 13.99$ & $53.88 \pm 22.90$ \\
		Trial 3 & Blues Solo & $85.12 \pm 21.50$ & $70.04 \pm 17.28$ & $58.54 \pm 19.01$ \\
		Trial 4 & Phase & $83.58 \pm 18.74$ & $66.88 \pm 18.94$ & $49.04 \pm 19.35$ \\
		Trial 5 & Ambient Guitar & $81.42 \pm 19.66$ & $74.15 \pm 15.57$ & $45.42 \pm 25.44$ \\
		\bottomrule
	\end{tabular}
\end{table}

\end{document}
